\documentclass[12pt, letterpaper]{article}

% ═══════════════════════════════════════════════════════════════
%  PAQUETES
% ═══════════════════════════════════════════════════════════════
\usepackage[utf8]{inputenc}
\usepackage[T1]{fontenc}
\usepackage[spanish]{babel}
\usepackage[top=2.5cm, bottom=2.5cm, left=3cm, right=3cm,
            headheight=1.25cm, footskip=1.25cm]{geometry}

% Colores
\usepackage[dvipsnames, table]{xcolor}
\definecolor{luluPurple}{HTML}{6B3FA0}
\definecolor{luluTeal}{HTML}{2DC7C7}
\definecolor{luluPink}{HTML}{E91E8C}
\definecolor{luluLila}{HTML}{C084FC}
\definecolor{luluDark}{HTML}{1A1A2E}
\definecolor{luluMidGray}{HTML}{888888}
\definecolor{luluLightGray}{HTML}{F0F0F0}
\definecolor{wireGray}{HTML}{CCCCCC}

% Tipografía
\usepackage{lmodern}
\usepackage{microtype}

% Interlineado y espaciado
\usepackage{setspace}
\setstretch{1.15}
\setlength{\parskip}{0pt plus 1pt}
\setlength{\parindent}{0pt}

% Gráficos y diagramas
\usepackage{graphicx}
\usepackage{tikz}
\usetikzlibrary{shapes, arrows.meta, positioning, fit, calc}

% Figuras
\usepackage{float}
\usepackage{caption}
\captionsetup{font=small, labelfont=bf, labelsep=period, skip=6pt}

% Tablas
\usepackage{tabularx}
\usepackage{booktabs}
\usepackage{array}
\newcolumntype{L}[1]{>{\raggedright\arraybackslash}p{#1}}

% Listas
\usepackage{enumitem}
\setlist[itemize]{leftmargin=*, topsep=3pt, itemsep=2pt, parsep=0pt}
\setlist[enumerate]{leftmargin=*, topsep=3pt, itemsep=2pt, parsep=0pt}
\setlist[description]{leftmargin=0pt, style=nextline, topsep=4pt, itemsep=6pt}

% Índice compacto (solo secciones principales)
\usepackage{tocloft}
\setlength{\cftbeforesecskip}{6pt}
\setcounter{tocdepth}{1}

% Formato de secciones
\usepackage{titlesec}
\titleformat{\section}
  {\large\bfseries\color{luluPurple}}{\thesection.}{0.6em}{}
  [\vspace{2pt}\color{luluPurple}\rule{\textwidth}{0.8pt}\vspace{4pt}]
\titleformat{\subsection}
  {\needspace{5\baselineskip}\normalsize\bfseries\color{luluDark}}{\thesubsection.}{0.5em}{}
\titlespacing{\section}{0pt}{14pt}{6pt}
\titlespacing{\subsection}{0pt}{8pt}{3pt}

% Encabezado y pie de página
\usepackage{fancyhdr}
\pagestyle{fancy}
\fancyhf{}
\fancyhead[L]{}
\fancyhead[R]{}
\setlength{\headheight}{14.2pt}
\fancyfoot[C]{\small\color{luluMidGray}\thepage}
\renewcommand{\headrulewidth}{0pt}
\renewcommand{\footrulewidth}{0pt}

% Cajas de color
\usepackage{tcolorbox}
\tcbuselibrary{skins, breakable}

% Hipervínculos
\usepackage[hidelinks]{hyperref}
\usepackage{url}

% Evitar que encabezados queden solos al final de página
\usepackage{needspace}


% ═══════════════════════════════════════════════════════════════
%  DOCUMENTO
% ═══════════════════════════════════════════════════════════════
\begin{document}


% ───────────────────────────────────────────────────────────────
%  PORTADA
% ───────────────────────────────────────────────────────────────
\begin{titlepage}
\pagecolor{luluDark}
\color{white}
\centering
\vspace*{2cm}

% Banda de marca
\begin{tikzpicture}
  \fill[luluPurple] (-7.5, -1.3) rectangle (7.5, 1.3);
  \node[font=\fontsize{28}{34}\selectfont\bfseries, white] at (0, 0.35) {LULÚ TAPIA};
  \node[font=\large\itshape, white!80] at (0, -0.55) {Arte que cobra vida.};
\end{tikzpicture}

\vspace{1.2cm}

{\LARGE\bfseries Propuesta de Sitio Web Personal}

\vspace{0.4cm}

{\large Diseño de Proyecto — Planeación y Desarrollo Web}

\vspace{2.2cm}

\begin{tabular}{rl}
  \textbf{Alumna:}             & Lourdes del Ángel Vásquez Tapia \\[7pt]
  \textbf{Plan de estudios:}   & Maestría en Diseño Multimedia \\[7pt]
  \textbf{Materia:}            & Arquitectura de información y diseño  \\[7pt]
  \textbf{Profesor(a):}        & María Isela García Velázquez       \\[7pt]
  \textbf{Institución:}        & Universidad Tecnológica de México (UNITEC)     \\[7pt]
  \textbf{Ciclo escolar:}      & Enero--Junio 2026 \\[7pt]
  \textbf{Fecha de entrega:}   & Junio de 2026                    \\
\end{tabular}

\vfill

\begin{tikzpicture}
  \fill[luluTeal] (-7.5, -0.25) rectangle (7.5, 0.25);
\end{tikzpicture}

\vspace{0.5cm}

\end{titlepage}

\pagecolor{white}
\color{black}

% Tabla de contenidos
{
  \setlength{\parskip}{0pt}
  \tableofcontents
}
\newpage


% ───────────────────────────────────────────────────────────────
%  RESUMEN EJECUTIVO
% ───────────────────────────────────────────────────────────────
\section*{Resumen Ejecutivo}
\addcontentsline{toc}{section}{Resumen Ejecutivo}

El presente documento constituye la propuesta integral de diseño y desarrollo del sitio web personal de \textbf{Lourdes del Ángel Vásquez Tapia}, artista freelance conocida profesionalmente como \textbf{Lulú Tapia}.

\textbf{Caso multimedia a solucionar:} Los artistas freelance del sector creativo digital enfrentan la problemática de no contar con un espacio propio que centralice su identidad de marca, portafolio y canal de contacto con clientes. La dependencia de plataformas externas (Behance, Instagram, ArtStation) fragmenta la presencia digital y limita el control sobre la experiencia del visitante y la captación de clientes potenciales.

\textbf{Metodología utilizada:} El proyecto adopta el marco de \textit{Diseño Centrado en el Usuario} (DCU) propuesto por Garrett~(2011), que estructura el desarrollo en cinco capas progresivas: estrategia, alcance, estructura, esqueleto y superficie visual. De forma complementaria, se integra el proceso de \textit{Design Thinking} de IDEO~(2015) para validar las necesidades reales del usuario en las etapas de empatía y definición.

\textbf{Contexto de desarrollo:} El proyecto se desarrolla en el ámbito académico de la Maestría en Diseño Multimedia de la UNITEC, con aplicación práctica real al contexto profesional de la artista. El sitio está destinado a operar en el mercado de servicios creativos digitales hispanohablante e internacional.

\textbf{Producto/solución:} Sitio web estático de una sola página con seis secciones interactivas (Inicio, Portafolio, Servicios, Sobre mí, Blog y Contacto), desarrollado en HTML5, CSS3 y JavaScript. El sitio puede consultarse en: \url{https://lulutapia.vercel.app}

\textbf{Características multimedia del diseño final:}
\begin{itemize}
  \item Interfaz web responsiva con tema visual oscuro (\texttt{\#1A1A2E}) y paleta cromática de alta saturación
  \item Galería de portafolio filtrable por categoría artística mediante JavaScript
  \item Animaciones de entrada al scroll con \texttt{IntersectionObserver}
  \item Efectos \textit{hover} interactivos sobre tarjetas de portafolio y servicios
  \item Formulario de contacto funcional con validación del lado del cliente
  \item Navegación responsiva con menú hamburguesa para dispositivos móviles

\end{itemize}

\textbf{Áreas de diseño multimedia favorecidas:} El proyecto beneficia principalmente las áreas de \textit{diseño de experiencia de usuario} (UX), \textit{diseño de interfaces} (UI), \textit{diseño gráfico digital} y \textit{comunicación visual}. Los hallazgos durante el desarrollo confirman que la jerarquía tipográfica y la paleta cromática cohesionada son determinantes en la percepción de profesionalismo de una marca personal creativa.

\newpage


% ───────────────────────────────────────────────────────────────
%  1. PLANTEAMIENTO DEL PROYECTO
% ───────────────────────────────────────────────────────────────
\section{Planteamiento del Proyecto del Sitio Web}

Este documento presenta la propuesta y planeación del sitio web personal de \textbf{Lourdes del Ángel Vásquez Tapia}, conocida profesionalmente como \textbf{Lulú Tapia}, artista freelance especializada en ilustración, modelado 3D, animación, storyboard, creación de personajes y vectorizado.

El sitio busca consolidar una identidad digital propia que sirva al mismo tiempo de portafolio artístico, escaparate de servicios y canal directo de contacto con clientes. El objetivo es posicionar a Lulú Tapia como marca personal reconocible dentro del mercado creativo digital, tanto a nivel nacional como internacional.

\subsection{Formato del Documento o Entregable}

Este entregable corresponde al \textbf{Diseño de Proyecto web}, presentado en formato \textbf{PDF}. Abarca los elementos de arquitectura de información: objetivos, definición de audiencia, modelo de negocio, mapa del sitio, sistema de navegación, estilo visual, sistema de organización y wireframes de las páginas principales. El documento se complementa con la implementación funcional del sitio web, accesible mediante la liga virtual incluida en la sección de Conclusión.

\subsection{Tentativo URL del Sitio}

\begin{tcolorbox}[
  colback=luluPurple!8, colframe=luluPurple,
  arc=4pt, boxrule=0.8pt,
  left=10pt, right=10pt, top=8pt, bottom=8pt
]
\centering
{\large\bfseries\color{luluPurple} www.lulutapia.com}\\[5pt]
{\small\color{luluMidGray} Alternativa: \textbf{www.lulutapia.art}}
\end{tcolorbox}

Se prefiere el dominio \texttt{.com} por su universalidad y reconocimiento global. El dominio \texttt{.art} se considera alternativa válida dado que fue creado específicamente para el sector creativo y artístico. En ambos casos, el nombre es corto, memorable y coincide directamente con el nombre de marca.

\subsection{Metodología de Diseño}

El desarrollo del sitio web de Lulú Tapia se fundamenta en la integración de dos metodologías ampliamente reconocidas en el campo del diseño de experiencias digitales:

\textbf{1. The Elements of User Experience — Jesse James Garrett (2011):} Garrett propone un modelo de cinco planos para el diseño de experiencias web: \textit{estrategia} (objetivos del sitio y necesidades del usuario), \textit{alcance} (funciones y contenido), \textit{estructura} (arquitectura de información), \textit{esqueleto} (diseño de interfaz y navegación) y \textit{superficie} (diseño visual). Este marco rige el presente documento: cada sección corresponde a un plano diferente de la metodología, desde los objetivos (Sección~2) hasta los wireframes (Sección~9) y el estilo visual (Sección~7). Garrett~(2011) subraya que ``la experiencia de usuario no ocurre por accidente; es el resultado de decisiones deliberadas y sistemáticas'' (p. 17), principio rector de este proyecto.

\textbf{2. Diseño Centrado en el Usuario — IDEO / Human-Centered Design (2015):} El proceso de IDEO estructura el diseño en tres fases: \textit{inspiración} (comprensión del contexto y las personas), \textit{ideación} (generación de soluciones) e \textit{implementación} (prototipado y entrega). La fase de inspiración se materializó en el análisis de audiencia (Sección~4); la ideación, en la definición de secciones, mapa del sitio y sistema de navegación (Secciones~3, 5 y 6); y la implementación, en los wireframes de alta y baja fidelidad y el desarrollo del sitio web funcional.

La integración de ambas metodologías permite un proceso riguroso: la metodología de Garrett garantiza la coherencia estructural del producto, mientras que el enfoque de IDEO asegura que las decisiones de diseño respondan a necesidades reales y verificadas del usuario final (Norman, 2013).


% ───────────────────────────────────────────────────────────────
%  2. TEMA, MENSAJE, METAS Y OBJETIVO
% ───────────────────────────────────────────────────────────────
\section{Tema, Mensaje, Metas y Objetivo}

\subsection{Tema}

El tema central del sitio es la \textbf{identidad artística y capacidad profesional de Lulú Tapia} como creadora multimedia. El sitio no es únicamente un portafolio, sino una experiencia visual que comunica el estilo, la versatilidad y la calidad del trabajo de la artista. El hilo conductor es el arte como puente entre la imaginación y la comunicación visual (Lupton, 2017).

\subsection{Mensaje Central}

\begin{tcolorbox}[
  colback=luluDark, colframe=luluDark,
  arc=4pt, boxrule=0pt,
  left=10pt, right=10pt, top=10pt, bottom=10pt
]
\centering
{\Large\itshape\color{white} ``Arte que cobra vida.''}
\end{tcolorbox}

Este mensaje resume la propuesta de valor de Lulú Tapia: su capacidad para llevar ideas, personajes e historias desde el concepto inicial hasta su realización en ilustración, movimiento y tridimensionalidad. Aplica tanto a clientes corporativos como a creadores independientes que buscan dar forma visual a sus proyectos.

\subsection{Metas}

\begin{enumerate}
  \item Establecer una presencia digital profesional y coherente con la identidad artística de la creadora.
  \item Atraer clientes potenciales de distintos sectores: entretenimiento, publicidad, videojuegos y medios digitales.
  \item Posicionarse como referente en ilustración y animación freelance dentro del mercado hispanohablante e internacional.
  \item Facilitar el contacto y la contratación directa de servicios creativos, eliminando intermediarios.
  \item Construir una comunidad de seguidores interesados en el proceso creativo y el desarrollo artístico de la marca.
\end{enumerate}

\subsection{Objetivo}

El sitio web funcionará como plataforma de presentación, captación de clientes y monetización de los servicios artísticos de Lulú Tapia; un punto de contacto directo entre la artista y quienes buscan sus servicios. Según Neumeier~(2006), ``una marca no es lo que tú dices que es; es lo que ellos dicen que es'' (p. 2), por lo que el diseño del sitio prioriza la percepción del cliente sobre la perspectiva interna.


% ───────────────────────────────────────────────────────────────
%  3. LISTA DE SECCIONES Y SUS FUNCIONES
% ───────────────────────────────────────────────────────────────
\section{Lista de Secciones y Sus Funciones del Sitio}

\begin{tcolorbox}[
  colback=luluLightGray, colframe=luluMidGray!40,
  arc=4pt, boxrule=0.6pt,
  left=8pt, right=8pt, top=6pt, bottom=6pt
]
El sitio se estructura en \textbf{seis secciones}, cada una con un rol definido dentro de la experiencia de navegación.
\end{tcolorbox}

\vspace{6pt}

\subsection{Inicio (Home)}
\textbf{Función:} Primera impresión visual del visitante. Contiene un encabezado impactante con el nombre de marca, tagline y una selección destacada del portafolio. Incluye llamados a la acción (CTA) directos hacia el portafolio y los servicios. Es la sección que define el tono visual y narrativo de todo el sitio, y el punto de entrada principal del embudo de conversión.

\subsection{Portafolio}
\textbf{Función:} Galería principal de trabajos organizados por categoría. Permite al visitante explorar el trabajo de Lulú Tapia por disciplina: ilustración, modelado 3D, storyboard, animación, creación de personajes y vectorizado. Cada proyecto puede incluir descripción, herramientas utilizadas y referencias del proceso creativo. Es la sección más importante para generar confianza y convertir visitantes en clientes. Nielsen y Loranger~(2006) señalan que la organización visual del contenido es el factor más relevante en la confianza inicial del usuario hacia un portafolio profesional.

\subsection{Servicios}
\textbf{Función:} Descripción clara y detallada de cada servicio ofrecido, incluyendo qué contempla cada uno, tiempos de entrega estimados y forma de solicitar una cotización. Funciona como carta de ventas directa y elimina la fricción en el proceso de contratación al responder las preguntas más comunes del cliente.

\subsection{Sobre mí}
\textbf{Función:} Presenta la historia personal y profesional de la artista, su trayectoria, filosofía de trabajo, herramientas que domina y valores. Humaniza la marca y genera confianza en el cliente potencial. Puede incluir fotografía o autorretrato, software que maneja y un curriculum vitae descargable.

\subsection{Proceso Creativo / Blog}
\textbf{Función:} Espacio donde Lulú Tapia comparte avances de proyectos, reflexiones sobre su proceso artístico y novedades. Posiciona a la artista como autoridad en su campo, mejora el posicionamiento SEO del sitio y fomenta el retorno de visitantes a través de contenido fresco y de valor.

\subsection{Contacto}
\textbf{Función:} Página dedicada a facilitar la comunicación directa. Incluye formulario de contacto, correo profesional, redes sociales activas y una guía breve sobre cómo iniciar un proyecto con la artista. Es el punto de cierre del embudo de conversión del sitio. Krug~(2014) establece que el camino hacia la acción debe ser intuitivo y requerir el mínimo de pasos posibles.


% ───────────────────────────────────────────────────────────────
%  4. AUDIENCIA, COMPETENCIAS Y MODELO DE NEGOCIO
% ───────────────────────────────────────────────────────────────
\section{Audiencia o Público y Competencias}

\subsection{Audiencia Principal}

El sitio está dirigido a cuatro segmentos de público objetivo:

\vspace{6pt}
\begin{tabularx}{\textwidth}{L{3cm} L{3.5cm} X}
\toprule
\rowcolor{luluPurple!10}
\textbf{Segmento} & \textbf{Perfil} & \textbf{Necesidad} \\
\midrule
Estudios de animación y producción & Directores creativos, productores & Artistas para proyectos de animación, storyboard o diseño de personajes \\
\addlinespace
Clientes independientes & Emprendedores, autores, creadores de contenido & Ilustraciones personalizadas, personajes para proyectos propios \\
\addlinespace
Desarrolladores de videojuegos & Indie devs y estudios pequeños & Assets visuales, concept art, sprites, modelado 3D \\
\addlinespace
Agencias y empresas & Directores de arte, marketers & Ilustración editorial, vectorizado, piezas para campañas publicitarias \\
\bottomrule
\end{tabularx}

\subsection{Competencias como Marca}

\begin{itemize}
  \item Dominio de múltiples disciplinas dentro del arte digital: ilustración, modelado 3D, animación, storyboard, creación de personajes y vectorizado.
  \item Capacidad de adaptación de estilo según el tipo de proyecto y las necesidades del cliente.
  \item Portafolio consolidado con proyectos variados y terminados listos para exhibir.
  \item Proceso de trabajo transparente y colaborativo con el cliente desde el concepto hasta la entrega.
\end{itemize}

\subsection{Modelo de Negocio}

El sitio generará ingresos a través de los siguientes mecanismos:

\begin{description}
  \item[\textbf{\color{luluPurple}1. Servicios freelance por proyecto}]
  Comisiones directas de ilustración, animación, storyboard, modelado 3D, creación de personajes y vectorizado, cotizadas según alcance y complejidad de cada encargo.

  \item[\textbf{\color{luluPurple}2. Paquetes de servicios}]
  Conjuntos predefinidos de entregables a precio fijo, por ejemplo: \textit{Paquete de personaje} (diseño + turnaround + hoja de expresiones) o \textit{Paquete editorial} (ilustración + vectorizado + variantes de color).

  \item[\textbf{\color{luluPurple}3. Licencias de uso de obra}]
  Venta de licencias comerciales para el uso de piezas ya creadas del portafolio, especialmente ilustraciones y diseños de personajes originales.

  \item[\textbf{\color{luluPurple}4. Contenido digital descargable} \normalfont\itshape (a futuro)]
  Venta de brushes, texturas, fondos, o personajes base que otros artistas o desarrolladores puedan adquirir y utilizar en sus propios proyectos.
\end{description}


% ───────────────────────────────────────────────────────────────
%  5. MAPA DEL SITIO
% ───────────────────────────────────────────────────────────────
\section{Mapa del Sitio}

El diagrama representa la arquitectura de información del sitio en tres niveles: dominio, secciones principales y subsecciones de contenido.

\vspace{12pt}

\begin{figure}[H]
\centering
\begin{tikzpicture}[
  every node/.style={
    draw, rounded corners=4pt,
    font=\small, align=center,
    inner sep=5pt
  },
  root/.style={
    fill=luluPurple, text=white,
    font=\small\bfseries,
    minimum width=4.5cm, minimum height=0.7cm
  },
  main/.style={
    fill=luluTeal!22, draw=luluTeal!70,
    font=\small\bfseries,
    minimum width=2.1cm, minimum height=0.65cm
  },
  sub/.style={
    fill=gray!10, draw=gray!40,
    font=\scriptsize,
    minimum width=2cm, minimum height=0.55cm
  },
  arr/.style={-{Stealth[length=3.5pt]}, gray!55, line width=0.7pt}
]

% Raíz
\node[root] (root) at (7, 11.5) {lulutapia.com};

% Nodos principales (nivel 1)
\node[main] (inicio)    at (1.1,  9.0) {Inicio};
\node[main] (port)      at (3.8,  9.0) {Portafolio};
\node[main] (serv)      at (6.5,  9.0) {Servicios};
\node[main] (sobre)     at (9.2,  9.0) {Sobre mí};
\node[main] (blog)      at (11.4, 9.0) {Blog};
\node[main] (cont)      at (13.3, 9.0) {Contacto};

% Flechas raíz → principales
\foreach \n in {inicio, port, serv, sobre, blog, cont}
  \draw[arr] (root) -- (\n);

% Subnodos: Inicio
\node[sub] (i1) at (1.1, 7.2) {Hero};
\node[sub] (i2) at (1.1, 6.4) {Destacados};
\node[sub] (i3) at (1.1, 5.6) {CTA};
\draw[arr] (inicio) -- (i1);
\draw[arr] (i1) -- (i2);
\draw[arr] (i2) -- (i3);

% Subnodos: Portafolio
\node[sub] (p1) at (3.8, 7.2) {Ilustración};
\node[sub] (p2) at (3.8, 6.4) {Modelado 3D};
\node[sub] (p3) at (3.8, 5.6) {Animación};
\node[sub] (p4) at (3.8, 4.8) {Personajes};
\node[sub] (p5) at (3.8, 4.0) {Storyboard};
\node[sub] (p6) at (3.8, 3.2) {Vectorizado};
\draw[arr] (port) -- (p1);
\draw[arr] (p1) -- (p2);
\draw[arr] (p2) -- (p3);
\draw[arr] (p3) -- (p4);
\draw[arr] (p4) -- (p5);
\draw[arr] (p5) -- (p6);

% Subnodos: Servicios
\node[sub] (s1) at (6.5, 7.2) {Lista de servicios};
\node[sub] (s2) at (6.5, 6.4) {Descripción};
\node[sub] (s3) at (6.5, 5.6) {Cotización};
\draw[arr] (serv) -- (s1);
\draw[arr] (s1) -- (s2);
\draw[arr] (s2) -- (s3);

% Subnodos: Sobre mí
\node[sub] (sm1) at (9.2, 7.2) {Biografía};
\node[sub] (sm2) at (9.2, 6.4) {Habilidades};
\node[sub] (sm3) at (9.2, 5.6) {CV descargable};
\draw[arr] (sobre) -- (sm1);
\draw[arr] (sm1) -- (sm2);
\draw[arr] (sm2) -- (sm3);

% Subnodos: Blog
\node[sub] (b1) at (11.4, 7.2) {Entradas};
\node[sub] (b2) at (11.4, 6.4) {Proceso};
\draw[arr] (blog) -- (b1);
\draw[arr] (b1) -- (b2);

% Subnodos: Contacto
\node[sub] (c1) at (13.3, 7.2) {Formulario};
\node[sub] (c2) at (13.3, 6.4) {Redes sociales};
\draw[arr] (cont) -- (c1);
\draw[arr] (c1) -- (c2);

\end{tikzpicture}
\caption{Mapa del sitio de lulutapia.com — arquitectura de información en tres niveles jerárquicos.}
\label{fig:mapa-sitio}
\end{figure}


% ───────────────────────────────────────────────────────────────
%  6. DISEÑO DE NAVEGACIÓN POR CONTENIDOS
% ───────────────────────────────────────────────────────────────
\section{Diseño de Navegación por Contenidos}

\subsection{Tipo de Navegación}

Menú de navegación horizontal fijo en la parte superior (\textit{sticky navigation}), visible en todo momento durante el desplazamiento de la página. Morville y Rosenfeld~(2006) señalan que la navegación persistente reduce la carga cognitiva del usuario al proporcionar un punto de referencia constante sobre la posición dentro del sitio.

\subsection{Estructura de Navegación Principal}

\begin{tcolorbox}[
  colback=luluDark, colframe=luluDark,
  arc=4pt, boxrule=0pt,
  left=8pt, right=8pt, top=8pt, bottom=8pt
]
\centering
{\small\bfseries\color{white}
Inicio \quad\textbar\quad Portafolio \quad\textbar\quad Servicios \quad\textbar\quad Sobre mí \quad\textbar\quad Blog \quad\textbar\quad Contacto}
\end{tcolorbox}

\subsection{Navegación Secundaria}

\begin{itemize}
  \item \textbf{En Portafolio:} sistema de filtros por categoría (botones de selección) que permite explorar el trabajo sin abandonar la página.
  \item \textbf{En Inicio:} botones CTA que dirigen al portafolio completo y a la sección de servicios.
  \item \textbf{En footer:} accesos directos a todas las secciones principales y enlaces a redes sociales activas.
\end{itemize}

\subsection{Navegación en Dispositivos Móviles}

Menú hamburguesa que despliega las opciones en formato vertical. El portafolio utilizará una cuadrícula adaptable (\textit{responsive grid}) de 1 a 3 columnas según el ancho de pantalla disponible.

\subsection{Flujo Principal del Usuario}

\vspace{8pt}
\begin{figure}[H]
\centering
\begin{tikzpicture}[
  node distance=1.2cm,
  box/.style={
    draw, rounded corners=3pt,
    fill=luluPurple!10, draw=luluPurple!50,
    font=\small, inner sep=7pt,
    minimum width=2.2cm, align=center
  },
  arrow/.style={-{Stealth[length=4pt]}, luluPurple!70, line width=0.9pt},
  label/.style={font=\scriptsize, text=luluMidGray}
]
\node[box] (a) {Inicio};
\node[box, right=of a] (b) {Portafolio};
\node[box, right=of b] (c) {Servicios};
\node[box, right=of c] (d) {Contacto};

\node[label, below=0.2cm of a] {Impacto visual};
\node[label, below=0.2cm of b] {Confianza};
\node[label, below=0.2cm of c] {Interés};
\node[label, below=0.2cm of d] {Conversión};

\draw[arrow] (a) -- (b);
\draw[arrow] (b) -- (c);
\draw[arrow] (c) -- (d);
\end{tikzpicture}
\caption{Flujo principal del usuario: embudo de conversión de cuatro etapas desde el primer contacto hasta la acción de contratar.}
\label{fig:flujo-usuario}
\end{figure}
\vspace{4pt}


% ───────────────────────────────────────────────────────────────
%  7. DESCRIPCIÓN DEL ESTILO VISUAL
% ───────────────────────────────────────────────────────────────
\section{Descripción del Estilo Visual}

El estilo visual es \textbf{colorido y expresivo}, alineado con la identidad artística de Lulú Tapia. El diseño parte de su propia obra: vivo, con carácter y fácil de recordar. Golombisky y Hagen~(2017) afirman que ``el color es la primera impresión que el ojo percibe y la más difícil de corregir una vez establecida'' (p. 112), razón por la cual la definición de paleta fue el primer paso de la fase de diseño.

\subsection{Paleta de Colores}

\vspace{10pt}
\begin{figure}[H]
\centering
\begin{tikzpicture}
  \fill[luluPurple,    rounded corners=4pt] (0.0,  0) rectangle (2.1, 1.9);
  \fill[luluTeal,      rounded corners=4pt] (2.3,  0) rectangle (4.4, 1.9);
  \fill[luluPink,      rounded corners=4pt] (4.6,  0) rectangle (6.7, 1.9);
  \fill[luluLila,      rounded corners=4pt] (6.9,  0) rectangle (9.0, 1.9);
  \fill[luluDark,      rounded corners=4pt] (9.2,  0) rectangle (11.3,1.9);
  \fill[luluLightGray, rounded corners=4pt] (11.5, 0) rectangle (13.6,1.9);
  \draw[gray!35,       rounded corners=4pt] (11.5, 0) rectangle (13.6,1.9);

  \node[white,    font=\small\bfseries]        at (1.05,  1.35) {Púrpura};
  \node[white,    font=\scriptsize\ttfamily]   at (1.05,  0.62) {\#6B3FA0};

  \node[white,    font=\small\bfseries]        at (3.35,  1.35) {Turquesa};
  \node[white,    font=\scriptsize\ttfamily]   at (3.35,  0.62) {\#2DC7C7};

  \node[white,    font=\small\bfseries]        at (5.65,  1.35) {Rosa};
  \node[white,    font=\scriptsize\ttfamily]   at (5.65,  0.62) {\#E91E8C};

  \node[white,    font=\small\bfseries]        at (7.95,  1.35) {Lila};
  \node[white,    font=\scriptsize\ttfamily]   at (7.95,  0.62) {\#C084FC};

  \node[white,    font=\small\bfseries]        at (10.25, 1.35) {Oscuro};
  \node[white,    font=\scriptsize\ttfamily]   at (10.25, 0.62) {\#1A1A2E};

  \node[black!60, font=\small\bfseries]        at (12.55, 1.35) {Claro};
  \node[black!50, font=\scriptsize\ttfamily]   at (12.55, 0.62) {\#F0F0F0};
\end{tikzpicture}
\caption{Paleta cromática de Lulú Tapia: seis tonos que definen la identidad visual del sitio.}
\label{fig:paleta}
\end{figure}
\vspace{6pt}

\subsection{Tipografía}

\begin{itemize}
  \item \textbf{Encabezados:} \textit{Baloo~2}, tipografía display de trazo amigable y expresivo, que aporta calidez y personalidad a los títulos sin sacrificar legibilidad.
  \item \textbf{Cuerpo de texto:} \textit{Inter}, tipografía sans-serif diseñada específicamente para pantallas, garantiza la lectura cómoda en todos los dispositivos y tamaños.
  \item El contraste entre titulares expresivos y texto funcional refuerza la jerarquía visual y guía la lectura (Lupton, 2017).
\end{itemize}

\subsection{Elementos Visuales}

\begin{itemize}
  \item Tema oscuro con fondo \texttt{\#1A1A2E} y tarjetas en \texttt{\#2A1845} / \texttt{\#221533}.
  \item Parches de color decorativos (rotados entre 5° y 15°) en tonos turquesa, rosa y lila distribuidos en fondos de sección.
  \item Confeti de puntos generado con gradientes CSS como textura de fondo sutil.
  \item Transiciones suaves y animaciones de entrada en los elementos al hacer scroll (\textit{fade-in}).
  \item Efectos \textit{hover} sobre las piezas del portafolio para revelar título y categoría del proyecto.
\end{itemize}

\subsection{Tono General}

Profesional pero cálido; creativo pero ordenado. El sitio debe sentirse como el estudio digital de un artista, no como un catálogo corporativo sin personalidad.


% ───────────────────────────────────────────────────────────────
%  8. SISTEMA DE ORGANIZACIÓN
% ───────────────────────────────────────────────────────────────
\section{Sistema de Organización}

El contenido del sitio se organiza de forma \textbf{jerárquica, con categorización por tipo de medio o servicio}.

\begin{description}
  \item[\textbf{\color{luluPurple}Nivel 1 — Navegación principal}]
  Las seis secciones del menú principal agrupan todo el contenido del sitio sin duplicidades. Cada sección tiene un propósito único y diferenciado dentro de la arquitectura de información (Morville y Rosenfeld, 2006).

  \item[\textbf{\color{luluPurple}Nivel 2 — Subsecciones}]
  Dentro del Portafolio, el contenido se organiza por disciplina artística (ilustración, 3D, animación, etc.). Dentro de Servicios, por tipo de entregable. Dentro del Blog, por fecha y categoría temática.

  \item[\textbf{\color{luluPurple}Etiquetado y filtrado}]
  Cada proyecto del portafolio lleva etiquetas de categoría, herramientas utilizadas y tipo de cliente, lo que facilita la búsqueda y el filtrado desde la interfaz sin necesidad de navegar a otra página.

  \item[\textbf{\color{luluPurple}Principio rector}]
  El usuario siempre debe poder ubicarse dentro del sitio y llegar a la sección de Contacto en \textbf{máximo tres clics} desde cualquier punto de la navegación (Krug, 2014).
\end{description}


% ───────────────────────────────────────────────────────────────
%  9. WIREFRAMES
% ───────────────────────────────────────────────────────────────
\section{Wireframes}

Los wireframes muestran la disposición de los elementos en las páginas principales, sin diseño visual definitivo ni paleta de color aplicada. Constituyen la capa de \textit{esqueleto} del modelo de Garrett~(2011), donde se define la distribución de componentes de interfaz y navegación antes de aplicar el diseño superficial.

% ┌──────────────────────────────────────────────────┐
%  WIREFRAME 1 — HOME
% └──────────────────────────────────────────────────┘
\subsection{Wireframe 1 — Inicio (Home)}

\begin{figure}[H]
\centering
\begin{tikzpicture}[scale=0.82, every node/.style={font=\scriptsize}]
  % Marco general
  \fill[gray!7] (0, 0) rectangle (14, 20);
  \draw[gray!40, line width=1pt] (0, 0) rectangle (14, 20);

  % Barra de dirección (simulada)
  \fill[gray!18] (0, 19.4) rectangle (14, 20);
  \node[gray!60] at (7, 19.7) {www.lulutapia.com};

  % Barra de navegación
  \fill[black!82] (0, 18.5) rectangle (14, 19.3);
  \node[white, anchor=west, font=\scriptsize\bfseries] at (0.4, 18.9) {Lulú Tapia};
  \node[white, font=\tiny] at (9.2, 18.9) {Inicio \;|\; Portafolio \;|\; Servicios \;|\; Sobre mí \;|\; Blog \;|\; Contacto};

  % Hero
  \fill[gray!12] (0, 13.2) rectangle (14, 18.4);
  \fill[gray!30] (0.4, 13.5) rectangle (6.8, 18.1);
  \node[gray!55] at (3.6, 15.8) {[ Imagen Hero ]};
  \node[anchor=west, font=\normalsize\bfseries] at (7.2, 17.4) {LULÚ TAPIA};
  \node[anchor=west, gray!65] at (7.2, 16.8) {Arte que cobra vida.};
  \fill[black!72] (7.2, 15.6) rectangle (10.5, 16.3);
  \node[white] at (8.85, 15.95) {Ver Portafolio};
  \draw[black!72, line width=0.8pt] (10.8, 15.6) rectangle (13.5, 16.3);
  \node[black!65] at (12.15, 15.95) {Contratar};

  % Sección: Trabajos destacados
  \node[anchor=west, font=\small\bfseries] at (0.4, 12.8) {Trabajos Destacados};
  \draw[gray!25, line width=0.5pt] (0.4, 12.55) -- (13.6, 12.55);

  \fill[gray!25] (0.4,  9.2) rectangle (3.8, 12.2);  \node[gray!55] at (2.1,  10.7) {Proyecto 1};
  \fill[gray!25] (4.1,  9.2) rectangle (7.1, 12.2);  \node[gray!55] at (5.6,  10.7) {Proyecto 2};
  \fill[gray!25] (7.4,  9.2) rectangle (10.4,12.2);  \node[gray!55] at (8.9,  10.7) {Proyecto 3};
  \fill[gray!25] (10.7, 9.2) rectangle (13.6,12.2);  \node[gray!55] at (12.15,10.7) {Proyecto 4};

  % CTA
  \fill[gray!16] (0, 7.3) rectangle (14, 8.9);
  \node[font=\small\bfseries] at (7, 8.45) {¿Tienes un proyecto en mente?};
  \fill[black!68] (5.0, 7.5) rectangle (9.0, 8.1);
  \node[white] at (7.0, 7.8) {Contáctame};

  % Footer
  \fill[black!82] (0, 0) rectangle (14, 3.5);
  \node[white, font=\small\bfseries, anchor=west] at (0.5, 3.1) {Lulú Tapia};
  \node[white, font=\tiny, anchor=west] at (0.5, 2.55) {Inicio \;|\; Portafolio \;|\; Servicios \;|\; Sobre mí \;|\; Blog \;|\; Contacto};
  \node[white, font=\tiny, anchor=west] at (0.5, 2.05) {Instagram \;|\; Behance \;|\; ArtStation \;|\; Pinterest};
  \node[white, font=\tiny, anchor=west] at (0.5, 1.5) {contacto@lulutapia.com};
  \draw[gray!40, line width=0.4pt] (0.4, 0.75) -- (13.6, 0.75);
  \node[gray!50, font=\tiny] at (7, 0.35) {© 2026 Lulú Tapia — Todos los derechos reservados};
\end{tikzpicture}
\caption{Wireframe 1 — Inicio (Home): disposición de hero, trabajos destacados, CTA y footer.}
\label{fig:wf-home}
\end{figure}

% ┌──────────────────────────────────────────────────┐
%  WIREFRAME 2 — PORTAFOLIO
% └──────────────────────────────────────────────────┘
\newpage
\subsection{Wireframe 2 — Portafolio}

\begin{figure}[H]
\centering
\begin{tikzpicture}[scale=0.82, every node/.style={font=\scriptsize}]
  \fill[gray!7] (0, 0) rectangle (14, 20);
  \draw[gray!40, line width=1pt] (0, 0) rectangle (14, 20);

  \fill[gray!18] (0, 19.4) rectangle (14, 20);
  \node[gray!60] at (7, 19.7) {www.lulutapia.com/portafolio};

  \fill[black!82] (0, 18.5) rectangle (14, 19.3);
  \node[white, anchor=west, font=\scriptsize\bfseries] at (0.4, 18.9) {Lulú Tapia};
  \node[white, font=\tiny] at (9.2, 18.9) {Inicio \;|\; \underline{Portafolio} \;|\; Servicios \;|\; Sobre mí \;|\; Blog \;|\; Contacto};

  \node[font=\large\bfseries] at (7, 17.9) {PORTAFOLIO};
  \node[gray!60] at (7, 17.4) {Explora mi trabajo por categoría};

  % Filtros — 7 botones
  \foreach \idx/\lbl in {0/Todos, 1/Ilustración, 2/3D, 3/Animación, 4/Personajes, 5/Storyboard, 6/Vectorizado} {
    \draw[gray!50, rounded corners=2pt]
      ({0.3 + \idx * 1.94}, 16.0) rectangle ({2.1 + \idx * 1.94}, 16.7);
    \node at ({1.2 + \idx * 1.94}, 16.35) {\lbl};
  }

  % Fila 1
  \fill[gray!25] (0.4, 12.0) rectangle (4.8, 15.6);  \node[gray!55] at (2.6, 13.8) {Ilustración 01};
  \fill[gray!25] (5.1, 12.0) rectangle (9.1, 15.6);  \node[gray!55] at (7.1, 13.8) {3D — Personaje};
  \fill[gray!25] (9.4, 12.0) rectangle (13.6,15.6);  \node[gray!55] at (11.5,13.8) {Animación 01};

  % Fila 2
  \fill[gray!25] (0.4,  7.4) rectangle (4.8, 11.6);  \node[gray!55] at (2.6,  9.5) {Storyboard 01};
  \fill[gray!25] (5.1,  7.4) rectangle (9.1, 11.6);  \node[gray!55] at (7.1,  9.5) {Personaje Sheet};
  \fill[gray!25] (9.4,  7.4) rectangle (13.6,11.6);  \node[gray!55] at (11.5, 9.5) {Vectorizado 01};

  \fill[black!82] (0, 0) rectangle (14, 3.5);
  \node[white, font=\small\bfseries, anchor=west] at (0.5, 3.1) {Lulú Tapia};
  \node[white, font=\tiny, anchor=west] at (0.5, 2.55) {Inicio \;|\; Portafolio \;|\; Servicios \;|\; Sobre mí \;|\; Blog \;|\; Contacto};
  \node[white, font=\tiny, anchor=west] at (0.5, 2.05) {Instagram \;|\; Behance \;|\; ArtStation \;|\; Pinterest};
  \node[white, font=\tiny, anchor=west] at (0.5, 1.5) {contacto@lulutapia.com};
  \draw[gray!40, line width=0.4pt] (0.4, 0.75) -- (13.6, 0.75);
  \node[gray!50, font=\tiny] at (7, 0.35) {© 2026 Lulú Tapia — Todos los derechos reservados};
\end{tikzpicture}
\caption{Wireframe 2 — Portafolio: galería filtrable por categoría artística con grid de 3 columnas.}
\label{fig:wf-portafolio}
\end{figure}

% ┌──────────────────────────────────────────────────┐
%  WIREFRAME 3 — SERVICIOS
% └──────────────────────────────────────────────────┘
\newpage
\subsection{Wireframe 3 — Servicios}

\begin{figure}[H]
\centering
\begin{tikzpicture}[scale=0.82, every node/.style={font=\scriptsize}]
  \fill[gray!7] (0, 0) rectangle (14, 20);
  \draw[gray!40, line width=1pt] (0, 0) rectangle (14, 20);

  \fill[gray!18] (0, 19.4) rectangle (14, 20);
  \node[gray!60] at (7, 19.7) {www.lulutapia.com/servicios};

  \fill[black!82] (0, 18.5) rectangle (14, 19.3);
  \node[white, anchor=west, font=\scriptsize\bfseries] at (0.4, 18.9) {Lulú Tapia};
  \node[white, font=\tiny] at (9.2, 18.9) {Inicio \;|\; Portafolio \;|\; \underline{Servicios} \;|\; Sobre mí \;|\; Blog \;|\; Contacto};

  \node[font=\large\bfseries] at (7, 17.9) {SERVICIOS};
  \node[gray!60] at (7, 17.4) {¿Qué puedo hacer por ti?};

  % Tarjeta: Ilustración
  \fill[gray!13] (0.4, 13.6) rectangle (4.6, 17.0);
  \fill[gray!28] (1.0, 15.7) rectangle (4.0, 16.7); \node[gray!55] at (2.5, 16.2) {[Ícono]};
  \node[font=\small\bfseries] at (2.5, 15.35) {Ilustración};
  \node[gray!60, font=\tiny] at (2.5, 14.8) {Ilustración digital a medida};
  \fill[black!68] (1.3, 13.8) rectangle (3.7, 14.3); \node[white] at (2.5, 14.05) {Cotizar};

  % Tarjeta: Modelado 3D
  \fill[gray!13] (5.0, 13.6) rectangle (9.0, 17.0);
  \fill[gray!28] (5.6, 15.7) rectangle (8.4, 16.7); \node[gray!55] at (7.0, 16.2) {[Ícono]};
  \node[font=\small\bfseries] at (7.0, 15.35) {Modelado 3D};
  \node[gray!60, font=\tiny] at (7.0, 14.8) {Modelado y render 3D};
  \fill[black!68] (5.9, 13.8) rectangle (8.1, 14.3); \node[white] at (7.0, 14.05) {Cotizar};

  % Tarjeta: Storyboard
  \fill[gray!13] (9.6, 13.6) rectangle (13.6,17.0);
  \fill[gray!28] (10.2,15.7) rectangle (13.0,16.7); \node[gray!55] at (11.6, 16.2) {[Ícono]};
  \node[font=\small\bfseries] at (11.6, 15.35) {Storyboard};
  \node[gray!60, font=\tiny] at (11.6, 14.8) {Guión gráfico para animación};
  \fill[black!68] (10.5,13.8) rectangle (12.7,14.3); \node[white] at (11.6, 14.05) {Cotizar};

  % Tarjeta: Animación
  \fill[gray!13] (0.4,  9.4) rectangle (4.6, 13.2);
  \fill[gray!28] (1.0, 11.5) rectangle (4.0, 12.9); \node[gray!55] at (2.5, 12.2) {[Ícono]};
  \node[font=\small\bfseries] at (2.5, 11.15) {Animación};
  \node[gray!60, font=\tiny] at (2.5, 10.55) {Animación 2D y motion graphics};
  \fill[black!68] (1.3,  9.6) rectangle (3.7, 10.1); \node[white] at (2.5, 9.85) {Cotizar};

  % Tarjeta: Creación de Personajes
  \fill[gray!13] (5.0,  9.4) rectangle (9.0, 13.2);
  \fill[gray!28] (5.6, 11.5) rectangle (8.4, 12.9); \node[gray!55] at (7.0, 12.2) {[Ícono]};
  \node[font=\small\bfseries, align=center] at (7.0, 11.1) {Creación de\\Personajes};
  \node[gray!60, font=\tiny] at (7.0, 10.55) {Diseño y hoja de personaje};
  \fill[black!68] (5.9,  9.6) rectangle (8.1, 10.1); \node[white] at (7.0, 9.85) {Cotizar};

  % Tarjeta: Vectorizado
  \fill[gray!13] (9.6,  9.4) rectangle (13.6,13.2);
  \fill[gray!28] (10.2,11.5) rectangle (13.0,12.9); \node[gray!55] at (11.6, 12.2) {[Ícono]};
  \node[font=\small\bfseries] at (11.6, 11.15) {Vectorizado};
  \node[gray!60, font=\tiny] at (11.6, 10.55) {Bocetos e imágenes a vector};
  \fill[black!68] (10.5, 9.6) rectangle (12.7,10.1); \node[white] at (11.6, 9.85) {Cotizar};

  \fill[black!82] (0, 0) rectangle (14, 3.5);
  \node[white, font=\small\bfseries, anchor=west] at (0.5, 3.1) {Lulú Tapia};
  \node[white, font=\tiny, anchor=west] at (0.5, 2.55) {Inicio \;|\; Portafolio \;|\; Servicios \;|\; Sobre mí \;|\; Blog \;|\; Contacto};
  \node[white, font=\tiny, anchor=west] at (0.5, 2.05) {Instagram \;|\; Behance \;|\; ArtStation \;|\; Pinterest};
  \node[white, font=\tiny, anchor=west] at (0.5, 1.5) {contacto@lulutapia.com};
  \draw[gray!40, line width=0.4pt] (0.4, 0.75) -- (13.6, 0.75);
  \node[gray!50, font=\tiny] at (7, 0.35) {© 2026 Lulú Tapia — Todos los derechos reservados};
\end{tikzpicture}
\caption{Wireframe 3 — Servicios: seis tarjetas de servicio con ícono, descripción y botón de cotización.}
\label{fig:wf-servicios}
\end{figure}

% ┌──────────────────────────────────────────────────┐
%  WIREFRAME 4 — CONTACTO
% └──────────────────────────────────────────────────┘
\newpage
\subsection{Wireframe 4 — Contacto}

\begin{figure}[H]
\centering
\begin{tikzpicture}[scale=0.82, every node/.style={font=\scriptsize}]
  \fill[gray!7] (0, 0) rectangle (14, 20);
  \draw[gray!40, line width=1pt] (0, 0) rectangle (14, 20);

  \fill[gray!18] (0, 19.4) rectangle (14, 20);
  \node[gray!60] at (7, 19.7) {www.lulutapia.com/contacto};

  \fill[black!82] (0, 18.5) rectangle (14, 19.3);
  \node[white, anchor=west, font=\scriptsize\bfseries] at (0.4, 18.9) {Lulú Tapia};
  \node[white, font=\tiny] at (9.2, 18.9) {Inicio \;|\; Portafolio \;|\; Servicios \;|\; Sobre mí \;|\; Blog \;|\; \underline{Contacto}};

  \node[font=\large\bfseries] at (7, 17.9) {CONTACTO};
  \node[gray!60] at (7, 17.4) {Hablemos de tu proyecto};

  % Columna izquierda: Formulario
  \node[anchor=west, font=\small\bfseries] at (0.4, 16.8) {Formulario de contacto};

  \node[anchor=west, gray!65] at (0.4, 16.00) {Nombre};
  \fill[white] (0.4, 15.15) rectangle (6.8, 15.65);
  \draw[gray!45, line width=0.5pt] (0.4, 15.15) rectangle (6.8, 15.65);

  \node[anchor=west, gray!65] at (0.4, 14.80) {Correo electrónico};
  \fill[white] (0.4, 13.95) rectangle (6.8, 14.45);
  \draw[gray!45, line width=0.5pt] (0.4, 13.95) rectangle (6.8, 14.45);

  \node[anchor=west, gray!65] at (0.4, 13.60) {Tipo de proyecto};
  \fill[white] (0.4, 12.75) rectangle (6.8, 13.25);
  \draw[gray!45, line width=0.5pt] (0.4, 12.75) rectangle (6.8, 13.25);

  \node[anchor=west, gray!65] at (0.4, 12.40) {Mensaje};
  \fill[white] (0.4, 10.05) rectangle (6.8, 12.05);
  \draw[gray!45, line width=0.5pt] (0.4, 10.05) rectangle (6.8, 12.05);

  \fill[black!75] (0.4, 9.05) rectangle (6.8, 9.70);
  \node[white, font=\small] at (3.6, 9.38) {Enviar mensaje};

  % Separador vertical
  \draw[gray!25, dashed, line width=0.6pt] (7.1, 8.85) -- (7.1, 16.9);

  % Columna derecha: Info de contacto
  \node[anchor=west, font=\small\bfseries] at (7.5, 16.8) {Información de contacto};
  \node[anchor=west, gray!65] at (7.5, 16.2) {contacto@lulutapia.com};
  \node[anchor=west, gray!65] at (7.5, 15.7) {Disponible para proyectos remotos};
  \node[anchor=west, gray!65] at (7.5, 15.2) {Respuesta en 24--48 horas};

  \node[anchor=west, font=\small\bfseries] at (7.5, 14.7) {Redes sociales};

  \foreach \sidx/\sname in {0/Instagram, 1/Behance, 2/ArtStation, 3/Pinterest} {
    \draw[gray!45, rounded corners=2pt]
      (7.5, {13.4 - \sidx * 0.85}) rectangle (13.5, {14.05 - \sidx * 0.85});
    \node[anchor=west] at (7.8, {13.725 - \sidx * 0.85}) {\sname};
  }

  % Footer
  \fill[black!82] (0, 0) rectangle (14, 3.5);
  \node[white, font=\small\bfseries, anchor=west] at (0.5, 3.1) {Lulú Tapia};
  \node[white, font=\tiny, anchor=west] at (0.5, 2.55) {Inicio \;|\; Portafolio \;|\; Servicios \;|\; Sobre mí \;|\; Blog \;|\; Contacto};
  \node[white, font=\tiny, anchor=west] at (0.5, 2.05) {Instagram \;|\; Behance \;|\; ArtStation \;|\; Pinterest};
  \node[white, font=\tiny, anchor=west] at (0.5, 1.5) {contacto@lulutapia.com};
  \draw[gray!40, line width=0.4pt] (0.4, 0.75) -- (13.6, 0.75);
  \node[gray!50, font=\tiny] at (7, 0.35) {© 2026 Lulú Tapia — Todos los derechos reservados};
\end{tikzpicture}
\caption{Wireframe 4 — Contacto: formulario de contacto y columna de información de redes sociales.}
\label{fig:wf-contacto}
\end{figure}


% ───────────────────────────────────────────────────────────────
%  10. PROCESO DE DISEÑO Y DESARROLLO
% ───────────────────────────────────────────────────────────────
\section{Proceso de Diseño y Desarrollo del Proyecto}

El desarrollo del sitio web de Lulú Tapia siguió un proceso estructurado de cinco etapas que integra las metodologías de Garrett~(2011) e IDEO~(2015) en acciones concretas y verificables. A continuación se describe cada etapa con sus actividades, decisiones clave y resultados obtenidos.

\subsection{Etapa 1 — Diagnóstico y Análisis de Necesidades}

El punto de partida fue identificar con claridad el problema a resolver: la artista contaba con trabajo de calidad pero carecía de una presencia digital centralizada que permitiera presentar su portafolio, comunicar sus servicios y recibir solicitudes de clientes de forma directa.

Para validar esta necesidad se realizó un análisis de audiencia (detallado en la Sección~4) que identificó cuatro grupos de clientes potenciales: estudios de animación y producción, clientes independientes, desarrolladores de videojuegos indie y agencias creativas. Asimismo, se exploró el panorama competitivo observando sitios de referencia en el mercado de ilustración y animación freelance, lo que permitió detectar patrones de diseño eficaces y oportunidades de diferenciación para la propuesta de Lulú Tapia.

Con base en este diagnóstico se definieron las seis secciones del sitio (Sección~3), el mensaje central —\textit{«Arte que cobra vida»}— y los objetivos de negocio (Sección~2), sentando así las bases estratégicas del proyecto.

\subsection{Etapa 2 — Arquitectura de Información y Estructura de Navegación}

Una vez definido el contenido necesario, se trabajó en la organización lógica del sitio. Se elaboró el mapa del sitio (Sección~5) representando de forma jerárquica la relación entre secciones y páginas, y se diseñó el flujo de navegación (Sección~6) para garantizar que el recorrido del usuario fuera intuitivo y condujera de manera natural desde el descubrimiento de la artista hasta el contacto directo.

Esta etapa corresponde a los planos de \textit{alcance} y \textit{estructura} del modelo de Garrett~(2011): se tomaron decisiones sobre qué contenido incluir y en qué orden lógico presentarlo, priorizando siempre la facilidad de uso sobre la complejidad visual. Se optó por una estructura de página única con desplazamiento vertical (\textit{one-page scroll}), ya que la naturaleza del sitio —portafolio personal— no requiere múltiples niveles de profundidad de contenido.

\subsection{Etapa 3 — Diseño Visual: Wireframes y Mockups de Alta Fidelidad}

Con la arquitectura definida, el proceso avanzó hacia la materialización visual del sitio. Esta etapa comprendió dos fases:

\textbf{Wireframes (baja fidelidad):} Se realizaron bocetos de las cuatro páginas principales —Inicio, Portafolio, Servicios y Contacto— sin diseño visual aplicado (Sección~9). El objetivo fue definir la disposición de los elementos de interfaz: encabezado, cuerpos de texto, galerías, formularios y llamadas a la acción, antes de introducir color o tipografía. Este paso redujo el riesgo de tomar decisiones visuales prematuras que pudieran comprometer la usabilidad.

\textbf{Mockups de alta fidelidad:} A partir de los wireframes se desarrollaron diseños visuales completos para las seis secciones del sitio —Inicio, Portafolio, Servicios, Sobre mí, Blog y Contacto— aplicando la paleta de color definida (\texttt{\#1A1A2E}, \texttt{\#6B3FA0}, \texttt{\#2DC7C7}, \texttt{\#E91E8C}), la tipografía seleccionada (\textit{Baloo 2} para títulos, \textit{Inter} para cuerpo) y la estética artesanal y expresiva que caracteriza la marca de la artista: fondo oscuro profundo, tarjetas con bordes iluminados, parches de color decorativos rotados y acentos vibrantes que refuerzan la identidad visual de Lulú Tapia. Los mockups sirvieron como guía visual definitiva para la etapa de implementación y como herramienta de validación con la cliente.

A continuación se presentan capturas representativas de los mockups de alta fidelidad correspondientes a las páginas de Inicio y Portafolio:

\begin{figure}[H]
\centering
\begin{tikzpicture}[scale=0.62]
  % Fondo de pagina
  \fill[luluDark] (0,0) rectangle (16,22);
  % Barra de navegacion
  \fill[black!60!luluDark] (0,21) rectangle (16,22);
  \fill[luluLila!70] (0.4,21.2) rectangle (4.2,21.8);
  \foreach \x in {6,8.2,10.4,12.6,14.3} {
    \fill[white,opacity=0.15] (\x,21.25) rectangle (\x+1.8,21.75);
  }
  % Hero — gradiente lateral manual con rectangulos
  \fill[luluPurple!35!luluDark] (0,15.5) rectangle (16,21);
  \fill[luluPurple!20!luluDark] (0,15.5) rectangle (10,21);
  \fill[luluDark] (0,15.5) rectangle (5,21);
  % Texto hero izquierda
  \fill[luluLila,opacity=0.9] (0.6,19.8) rectangle (7.5,20.5);
  \fill[white,opacity=0.3] (0.6,18.9) rectangle (5.8,19.3);
  \fill[white,opacity=0.2] (0.6,18.2) rectangle (5,18.55);
  \fill[luluPink,opacity=0.85] (0.6,16.8) rectangle (3.2,17.5);
  \fill[luluPurple,opacity=0.85] (3.5,16.8) rectangle (6.1,17.5);
  % Artwork derecha — dos imagenes superpuestas (representacion)
  \fill[luluPurple!30] (8.5,16) rectangle (15.4,20.8);
  \fill[luluPink!15] (8.5,16) rectangle (15.4,17.3);
  \fill[luluPurple!18] (9.2,16.5) rectangle (14.6,20.5);
  \draw[luluLila,opacity=0.35,line width=0.6pt] (8.5,16) rectangle (15.4,20.8);
  \node[white,opacity=0.45,font=\footnotesize] at (11.9,18.4) {Obra destacada};
  % Seccion portafolio
  \fill[black!35!luluDark] (0,8) rectangle (16,15.2);
  \fill[luluLila,opacity=0.7] (0.5,14.5) rectangle (4,15);
  % Filtros
  \foreach \x in {0.5,3.2,5.9,8.6,11.3} {
    \fill[luluPurple!45] (\x,13.6) rectangle (\x+2.4,14.1);
  }
  % Grid de piezas 3 columnas x 2 filas
  \foreach \x in {0.4,5.5,10.6} {
    \fill[luluPurple!22] (\x,8.3) rectangle (\x+4.8,11);
    \fill[luluPink!12] (\x,8.3) rectangle (\x+4.8,8.85);
    \fill[luluPurple!22] (\x,11.3) rectangle (\x+4.8,13.3);
    \fill[luluPink!12] (\x,11.3) rectangle (\x+4.8,11.85);
  }
  % Seccion servicios
  \fill[black!20!luluDark] (0,1.5) rectangle (16,7.7);
  \fill[luluLila,opacity=0.7] (0.5,7.2) rectangle (3.5,7.65);
  \foreach \x in {0.4,5.6,10.8} {
    \fill[black!30] (\x,1.8) rectangle (\x+4.8,6.9);
    \fill[luluTeal!40] (\x+0.3,5.8) rectangle (\x+2.2,6.5);
    \fill[white,opacity=0.12] (\x+0.3,3) rectangle (\x+4.3,5.5);
    \fill[luluPink!60] (\x+0.3,2) rectangle (\x+4.3,2.65);
  }
  % Footer
  \fill[black!75!luluDark] (0,0) rectangle (16,1.3);
\end{tikzpicture}
\caption{Mockup de alta fidelidad — Página de Inicio: hero con logo y tagline (izquierda), obra destacada (derecha), galería filtrable y sección de servicios.}
\label{fig:mockup-home}
\end{figure}

\begin{figure}[H]
\centering
\begin{tikzpicture}[scale=0.62]
  % Fondo de pagina
  \fill[luluDark] (0,0) rectangle (16,18);
  % Barra de navegacion
  \fill[black!60!luluDark] (0,17) rectangle (16,18);
  \fill[luluLila!70] (0.4,17.2) rectangle (4.2,17.8);
  \foreach \x in {6,8.2,10.4,12.6,14.3} {
    \fill[white,opacity=0.15] (\x,17.25) rectangle (\x+1.8,17.75);
  }
  % Encabezado seccion
  \fill[luluDark] (0,14.5) rectangle (16,17);
  \fill[luluLila,opacity=0.9] (0.5,16.2) rectangle (5,16.8);
  \fill[white,opacity=0.2] (0.5,15.4) rectangle (7,15.8);
  % Barra de filtros
  \fill[black!30!luluDark] (0,13.2) rectangle (16,14.2);
  \foreach \x in {0.5,3.5,6.5,9.5,12.5} {
    \fill[luluPurple!55] (\x,13.35) rectangle (\x+2.6,13.95);
  }
  % Grid portafolio 4 columnas x 3 filas
  \fill[black!25!luluDark] (0,0) rectangle (16,13);
  \foreach \x in {0.3,4.3,8.3,12.3} {
    \foreach \y in {0.4,4.5,8.6} {
      \fill[black!40] (\x,\y) rectangle (\x+3.8,\y+3.7);
      % Imagen placeholder con gradiente simulado
      \fill[luluPurple!25] (\x,\y) rectangle (\x+3.8,\y+2.8);
      \fill[luluPink!10] (\x,\y) rectangle (\x+3.8,\y+0.6);
      % Overlay info en hover
      \fill[luluPurple!60,opacity=0.15] (\x,\y+2.8) rectangle (\x+3.8,\y+3.7);
      \fill[white,opacity=0.12] (\x+0.2,\y+3.1) rectangle (\x+2.5,\y+3.4);
      \fill[luluTeal,opacity=0.5] (\x+0.2,\y+2.95) rectangle (\x+1.6,\y+3.05);
    }
  }
\end{tikzpicture}
\caption{Mockup de alta fidelidad — Portafolio: galería filtrable por categoría (ilustración, personajes, modelado 3D, storyboard) con vista previa de obras y datos al pasar el cursor.}
\label{fig:mockup-portafolio}
\end{figure}

\subsection{Etapa 4 — Implementación del Sitio Web}

La implementación tradujo los diseños aprobados en un sitio web funcional. Se eligió un desarrollo en tecnologías web fundamentales —HTML5 para la estructura del contenido, CSS3 para los estilos visuales y JavaScript para la interactividad— sin recurrir a dependencias de terceros, lo que simplifica el mantenimiento a largo plazo y mantiene el control total del código en manos del equipo de desarrollo.

Las funcionalidades clave implementadas fueron:

\begin{itemize}
  \item \textbf{Galería de portafolio filtrable:} Permite al visitante explorar las piezas por categoría (ilustración, personajes, modelado 3D, storyboard) con transiciones animadas al aplicar los filtros.
  \item \textbf{Vista ampliada de piezas (lightbox):} Al hacer clic sobre cualquier obra del portafolio, se despliega una vista de pantalla completa con el título y categoría de la pieza, sin abandonar la página.
  \item \textbf{Animaciones de aparición al desplazarse:} Los contenidos de cada sección aparecen con una transición suave conforme el visitante avanza hacia abajo en la página, lo que refuerza la percepción de dinamismo y cuidado visual.
  \item \textbf{Diseño adaptable (\textit{responsive}):} El sitio se adapta automáticamente a distintos tamaños de pantalla —computadoras de escritorio, tabletas y teléfonos móviles— garantizando una experiencia consistente en todos los dispositivos.
  \item \textbf{Formulario de contacto:} Permite a los clientes potenciales iniciar una solicitud directamente desde el sitio, con validación de campos y mensaje de confirmación de envío.
\end{itemize}

\subsection{Etapa 5 — Publicación y Verificación}

Una vez completada la implementación y realizadas pruebas de funcionamiento en distintos dispositivos y navegadores, el sitio fue publicado en línea mediante un servicio de hospedaje web gratuito, quedando accesible de forma inmediata al público general sin costo inicial para la artista.

Se verificó el correcto funcionamiento de todas las secciones, filtros, animaciones y formulario de contacto. El resultado es un sitio web profesional, de carga rápida y totalmente operativo, que cumple con los objetivos comunicacionales y comerciales definidos en la etapa de diagnóstico. La liga al sitio publicado se incluye en la conclusión del presente documento.

% ───────────────────────────────────────────────────────────────
%  11. CONCLUSIÓN
% ───────────────────────────────────────────────────────────────
\section{Conclusión}

El sitio web de Lulú Tapia parte de una necesidad concreta y verificada: los artistas freelance requieren un espacio digital propio desde el que presentarse, mostrar su trabajo y generar oportunidades de negocio sin depender de plataformas externas. La aplicación combinada de la metodología de Garrett~(2011) y el proceso de IDEO~(2015) proporcionó un marco de trabajo riguroso que garantiza tanto la coherencia estructural como la pertinencia hacia las necesidades del usuario.

La planeación aquí desarrollada define una estructura funcional y visualmente coherente con la identidad artística de Lourdes del Ángel Vásquez Tapia. Cada sección tiene un propósito claro dentro del recorrido del usuario, y el conjunto del sitio atiende tanto objetivos de comunicación como de negocio.

\textbf{Pertinencia y viabilidad:} La propuesta es pertinente porque responde a una demanda real del mercado creativo digital y se alinea con las prácticas documentadas de branding personal (Neumeier, 2006) y diseño de experiencia de usuario (Norman, 2013). Es viable técnicamente porque se implementa con tecnologías web estándar (HTML5, CSS3, JavaScript) sin dependencias externas, lo que facilita el mantenimiento y la actualización futura del sitio. La opción de hospedaje gratuito en Vercel hace accesible el despliegue sin costo inicial.

\textbf{Limitaciones:} El sitio en su estado actual no incluye un CMS (\textit{Content Management System}), por lo que la actualización del portafolio requiere editar el código fuente directamente. En una fase posterior se podría integrar una solución headless como Contentful o Sanity para separar contenido y código.

\textbf{Liga virtual al sitio implementado:} El sitio web funcional resultado de este proyecto puede consultarse en la siguiente dirección:

\begin{tcolorbox}[
  colback=luluPurple!8, colframe=luluPurple,
  arc=4pt, boxrule=0.8pt,
  left=10pt, right=10pt, top=8pt, bottom=8pt
]
\centering
{\large\bfseries\color{luluPurple} \url{https://lulutapia.vercel.app}}
\end{tcolorbox}


% ───────────────────────────────────────────────────────────────
%  11. REFERENCIAS BIBLIOGRÁFICAS
% ───────────────────────────────────────────────────────────────
\section{Referencias Bibliográficas}

\begin{thebibliography}{99}

\bibitem{clark2020}
Clark, B. (2020). \textit{The content marketer's guide to personal branding}. Copyblogger. \url{https://copyblogger.com}

\bibitem{garrett2011}
Garrett, J. J. (2011). \textit{The elements of user experience: User-centered design for the web and beyond} (2a ed.). New Riders.

\bibitem{golombisky2017}
Golombisky, K., y Hagen, R. (2017). \textit{White space is not your enemy: A beginner's guide to communicating visually through graphic, web \& multimedia design} (3a ed.). Focal Press.

\bibitem{ideo2015}
IDEO. (2015). \textit{The field guide to human-centered design}. IDEO.org. \url{https://www.designkit.org}

\bibitem{krug2014}
Krug, S. (2014). \textit{Don't make me think, revisited: A common sense approach to web usability} (3a ed.). New Riders.

\bibitem{lupton2017}
Lupton, E. (2017). \textit{Design is storytelling}. Cooper Hewitt, Smithsonian Design Museum.

\bibitem{morville2006}
Morville, P., y Rosenfeld, L. (2006). \textit{Information architecture for the world wide web} (3a ed.). O'Reilly Media.

\bibitem{neumeier2006}
Neumeier, M. (2006). \textit{The brand gap: How to bridge the distance between business strategy and design}. New Riders.

\bibitem{nielsen2006}
Nielsen, J., y Loranger, H. (2006). \textit{Prioritizing web usability}. New Riders.

\bibitem{norman2013}
Norman, D. A. (2013). \textit{The design of everyday things} (Ed. rev. y ampliada). Basic Books.

\end{thebibliography}

\end{document}
